\section{Discussion}
\label{sec:discussion}

\paragraph{Why the leakage was overlooked.}
Two factors. First, the leakage manifests as ``great training error and
plausible test error'', not the classic ``zero training error and broken
generalization'' pattern. Per-split z-score normalization scales the leak
to a different magnitude on test than on train, leaving a visible
generalization gap that masks the underlying issue. Second, the
direction of the leak (subtracting the target rather than adding it,
followed by normalization) means the network must do nontrivial work
to recover the target, lending plausibility to the learning curve. Our
audit took several weeks to isolate after we realized the published
numbers could not be reproduced from the released code on a clean
subject-disjoint split.

\paragraph{What this means for ``neural calibration''.}
Our negative results do not rule out neural refinement of eye-tracking
calibration in general. They do show that simply ensembling the
predictions of several classical calibrators and feeding them into an MLP
is insufficient on the spatial scales and noise regimes of head-mounted
trackers. We hypothesize that real wins will require either richer
inputs (raw per-fixation samples, pupil size, head pose) or richer
supervision (smooth-pursuit trajectories, naturally-occurring
saccades). Both directions are out of scope for the present submission.

\paragraph{Limitations of online shrinkage.}
The largest practical cost of our method is the requirement of
$K \approx 12$ extra fixations after the standard calibration phase --- on
the order of $30$ extra seconds of user time. The largest scientific
limitation is dataset scale: only three of our twelve subjects have
enough trials to admit $K=12$ LOSO episodes. We pre-registered our
intent to validate on a larger pool but could not collect new data
within the rebuttal window; the camera-ready version will report
results on twelve additional participants we are currently recording.

\paragraph{Why does the learned shrinkage outperform the global $\lambda$
only at $K \ge 12$?}
At small $K$ the per-trial bias statistics ($\bar{b}_K$, $S_K$) are
themselves noisy, so the network has limited signal to differentiate
trials. The fixed $\lambda^\star(K)$ is therefore very close to optimal.
At larger $K$ the bias statistics become stable, and the shrinkage
network can identify high-bias trials whose mean estimates are
trustworthy and downshrink low-bias trials whose mean estimates would
otherwise destabilize the anchor.

\paragraph{Honest reproducibility checklist.}
We release: (a) the original leaky implementation alongside the honest
patched version (a one-line diff to \texttt{model.py}); (b) the
target-disjoint JuDo1000 split with baselines fit on training targets
only; (c) the LOSO splits and the shrinkage MLP training script; (d)
JSON results for every experiment in this paper; (e) the dataset of
2\,526 fixations from 143 trials with per-fixation raw samples,
target labels, and pre-computed classical baseline predictions. All
artifacts are anonymous-friendly; subject identifiers are random
codes.

\paragraph{Conclusion.}
Eye tracking calibration is a useful test bed for honest evaluation of
neural augmentation. We found that a recently-proposed multi-baseline
neural refinement design owes its headline numbers to a single
input-feature leak, and that no honest variant of the design beats the
strongest classical baseline. We then proposed a $\sim$1\,200-parameter
shrinkage estimator that takes advantage of $K=12$ extra calibration
fixations to deliver a $17\%$ honest improvement over the strongest
classical baseline, evaluated leave-one-subject out. The result is
modest in absolute terms but reproducible, principled, and immediately
deployable. We hope the audit and the open release help establish a
firmer benchmarking culture in the calibration sub-field.

\begin{ack}
We thank the anonymous reviewers of an earlier version of this work
whose careful reading prompted our re-audit and ultimately led to the
discovery of the leakage.
\end{ack}
